%!TEX root = ./template-skripsi.tex
%-------------------------------------------------------------------------------
%                            	BAB V
%               		KESIMPULAN DAN SARAN
%-------------------------------------------------------------------------------

\chapter{KESIMPULAN DAN SARAN}

\section{Kesimpulan}
Berdasarkan hasil implementasi dan pengujian fitur sistem informasi yang telah dirancang, maka diperoleh kesimpulan sebagai berikut:

\begin{enumerate}
	\item Pengimplementasian algorithma penghitungan panjang dan berat ikan menggunakan metode \textit{image processing} dengan algoritma \textit{Harris corner detection} masih belum memberikan hasil yang akurat, hal ini disebabkan oleh faktor-faktor seperti kualitas citra, jarak kamera, sudut pengambilan gambar, posisi ikan yang tidak selalu ideal, dan bentuk tubuh ikan.
	\item Berdasarkan hasil evaluasi akurasi, rata-rata pada perhitungan panjang ikan mencapai 100\%, sedangkan pada perhitungan berat ikan rata-rata mencapai 74.16\%. Hal ini menunjukkan bahwa algorithma masih memiliki keterbatasan dalam mengestimasi berat ikan secara akurat, terutama untuk jenis ikan yang memiliki bentuk tubuh yang lebih kompleks.
	% \item Kurang-nya dataset yang digunakan untuk pengujian juga menjadi faktor 
\end{enumerate}

\section{Saran}
Adapun saran untuk penelitian selanjutnya adalah:
\begin{enumerate} 
	\item Menggunakan metode \textit{image processing} yang lebih canggih seperti \textit{deep learning} untuk meningkatkan akurasi penghitungan panjang dan berat ikan.
	\item Melakukan pengujian dengan dataset yang lebih besar dan beragam untuk memastikan keandalan sistem dalam berbagai kondisi.
	\item Mengoptimalkan algoritma yang digunakan untuk mempercepat proses penghitungan dan meningkatkan efisiensi sistem.
	\item Melakukan analisis lebih lanjut terhadap faktor-faktor yang mempengaruhi akurasi penghitungan, seperti pencahayaan, kualitas citra, dan posisi ikan, untuk mengembangkan solusi yang lebih robust.
\end{enumerate}


% Baris ini digunakan untuk membantu dalam melakukan sitasi
% Karena diapit dengan comment, maka baris ini akan diabaikan
% oleh compiler LaTeX.
\begin{comment}
\bibliography{daftar-pustaka}
\end{comment}
